\section{习题课 \#9（11月8日）}

\subsection{线性映射的核、像与矩阵表示}

设 $V_1,V_2$ 是域 $\F$ 上的向量空间。映射 $T:V_1\to V_2$ 称为线性映射，如果
\begin{equation}
 T(au+bv)=aT(u)+bT(v),
 \qquad a,b\in\F,
 \quad u,v\in V_1.
\end{equation}
其核与像分别为
\begin{equation}
 \Ker T=\set{v\in V_1:T(v)=0},
 \qquad
 \Image T=\set{T(v):v\in V_1}.
\end{equation}
$T$ 单射当且仅当 $\Ker T=\{0\}$，满射当且仅当 $\Image T=V_2$。

固定 $V_1,V_2$ 的基后，线性映射由一个矩阵表示。若更换定义域与值域的基，表示矩阵在左右分别乘以可逆矩阵，因此同一线性映射的不同矩阵表示彼此等价。

\begin{example}
令
\begin{equation}
 A=\begin{pmatrix}
 3&1&4&7\\
 -1&1&0&3
 \end{pmatrix}.
\end{equation}
前两列线性无关，后两列可由它们表示，所以
\begin{equation}
 \Image A
 =\Span\left\{
 \begin{pmatrix}3\\-1\end{pmatrix},
 \begin{pmatrix}1\\1\end{pmatrix}
 \right\},
 \qquad
 \rank A=2.
\end{equation}
由 $\dim\Ker A=4-2=2$，消元可得
\begin{equation}
 \Ker A=\Span\left\{
 \begin{pmatrix}-1\\-1\\1\\0\end{pmatrix},
 \begin{pmatrix}-1\\-4\\0\\1\end{pmatrix}
 \right\}.
\end{equation}
\end{example}

\subsection{矩阵方程 AXB=C}

先考虑可逆情形。若 $A,B$ 都可逆，则
\begin{equation}
 AXB=C
 \quad\Longleftrightarrow\quad
 X=A^{-1}CB^{-1}.
\end{equation}
一般情形可使用广义逆统一描述。

\begin{definition}
矩阵 $A^{-}$ 称为 $A$ 的一个内逆，如果
\begin{equation}
 AA^{-}A=A.
\end{equation}
每个有限维矩阵都存在内逆。事实上，若
\begin{equation}
 A=P\begin{pmatrix}I_r&0\\0&0\end{pmatrix}Q
\end{equation}
是秩标准形，则可取
\begin{equation}
 A^{-}=Q^{-1}
 \begin{pmatrix}I_r&0\\0&0\end{pmatrix}
 P^{-1},
\end{equation}
其中零块按尺寸补齐。
\end{definition}

\begin{theorem}[矩阵方程的相容条件与通解]
设
\begin{equation}
 A\in\F^{m\times n},
 \quad X\in\F^{n\times p},
 \quad B\in\F^{p\times q},
 \quad C\in\F^{m\times q},
\end{equation}
并任选内逆 $A^{-},B^{-}$。方程
\begin{equation}
 AXB=C
\end{equation}
有解当且仅当
\begin{equation}
 AA^{-}CB^{-}B=C.
\end{equation}
相容时，全部解为
\begin{equation}
 X=A^{-}CB^{-}+Y-A^{-}AYBB^{-},
 \qquad Y\in\F^{n\times p}\text{ 任意}.
\end{equation}
\end{theorem}

\begin{proof}
若 $AXB=C$，则
\begin{equation}
 AA^{-}CB^{-}B
 =AA^{-}AXBB^{-}B
 =AXB=C.
\end{equation}
反之，相容条件成立时，取 $X_0=A^{-}CB^{-}$ 即有 $AX_0B=C$。

对任意 $Y$，代入通解表达式得
\begin{align}
 AXB
 &=AA^{-}CB^{-}B+AYB-AA^{-}AYBB^{-}B\\
 &=C+AYB-AYB=C.
\end{align}
若 $X$ 是任一解，在右端取 $Y=X$，则
\begin{equation}
 A^{-}CB^{-}+X-A^{-}AXBB^{-}=X,
\end{equation}
故该公式确实包含全部解。
\end{proof}

\begin{remark}
相容条件可解释为两个空间条件：$C$ 的每一列必须属于 $\Image A$，而 $C$ 的每一行必须属于 $B$ 的行空间。若只求 $AX=C$，则对 $C$ 的每一列分别解同一个系数矩阵的方程即可；若只求 $XB=C$，可转置后同样处理。
\end{remark}

\subsection{秩标准形中的像与核}

设 $\rank A=r$，并写成
\begin{equation}
 A=P
 \begin{pmatrix}I_r&0\\0&0\end{pmatrix}
 Q,
 \qquad P,Q\text{ 可逆}.
\end{equation}
把
\begin{equation}
 P=(P_1,P_2),
 \qquad
 Q^{-1}=(Q_1,Q_2)
\end{equation}
按秩分块，则
\begin{equation}
 \Image A=\Image P_1,
 \qquad
 \Ker A=\Image Q_2.
\end{equation}
相应地，
\begin{equation}
 \Image A^T=\Image Q^T_1,
 \qquad
 \Ker A^T=\Image P^{-T}_2.
\end{equation}
这些公式把像空间、零空间和秩标准形直接联系起来。

\begin{proposition}
若 $A=P\begin{psmallmatrix}I_r&0\\0&0\end{psmallmatrix}Q$，则
\begin{equation}
 \Ker A=\Ker\begin{pmatrix}I_r&0\\0&0\end{pmatrix}Q,
\end{equation}
而左乘可逆矩阵不改变核，右乘可逆矩阵只作坐标变换。类似地，右乘可逆矩阵不改变列空间的维数，左乘可逆矩阵把列空间同构地变换。
\end{proposition}

\subsection{分块矩阵与 Schur 补}

\begin{theorem}
设
\begin{equation}
 M=\begin{pmatrix}A_{11}&A_{12}\\A_{21}&A_{22}\end{pmatrix}
\end{equation}
且 $A_{11}$ 可逆。记
\begin{equation}
 S=A_{22}-A_{21}A_{11}^{-1}A_{12}.
\end{equation}
则 $M$ 可逆当且仅当 $S$ 可逆，并且
\begin{equation}
 M^{-1}=\begin{pmatrix}
 A_{11}^{-1}+A_{11}^{-1}A_{12}S^{-1}A_{21}A_{11}^{-1}
 &-A_{11}^{-1}A_{12}S^{-1}\\
 -S^{-1}A_{21}A_{11}^{-1}&S^{-1}
 \end{pmatrix}.
\end{equation}
\end{theorem}

\begin{proof}
利用分解
\begin{equation}
 M=
 \begin{pmatrix}I&0\\A_{21}A_{11}^{-1}&I\end{pmatrix}
 \begin{pmatrix}A_{11}&0\\0&S\end{pmatrix}
 \begin{pmatrix}I&A_{11}^{-1}A_{12}\\0&I\end{pmatrix}
\end{equation}
即可。
\end{proof}

\subsection{若干秩恒等式与可解性判别}

\begin{theorem}
设 $A\in\F^{m\times n}$、$B\in\F^{n\times m}$。则
\begin{equation}
 \rank(A-ABA)=\rank A+\rank(I_n-BA)-n.
\end{equation}
\end{theorem}

\begin{proof}
写成
\begin{equation}
 A-ABA=A(I_n-BA).
\end{equation}
Sylvester 秩不等式给出下界。另一方面，若 $x\in\Ker A$，则
\begin{equation}
 (I_n-BA)x=x,
\end{equation}
所以 $\Ker A\subseteq\Image(I_n-BA)$。这正是 Sylvester 不等式取等号的条件。
\end{proof}

\begin{proposition}
设矩阵尺寸相容。方程
\begin{equation}
 ABX=A
\end{equation}
有解当且仅当
\begin{equation}
 \rank(AB)=\rank A.
\end{equation}
\end{proposition}

\begin{proof}
若有解，则 $\Image A\subseteq\Image(AB)$，而反向包含总成立，故两列空间相同。反之，若秩相同，则 $\Image A=\Image(AB)$，逐列求解即可构造 $X$。
\end{proof}

\begin{proposition}
对同阶方阵 $A,B$，
\begin{equation}
 \rank(I-AB)\leq\rank(I-A)+\rank(I-B).
\end{equation}
\end{proposition}

\begin{proof}
由
\begin{equation}
 I-AB=(I-A)+A(I-B)
\end{equation}
及矩阵和、矩阵乘积的秩上界立即得到。
\end{proof}

\begin{theorem}
若同阶方阵满足
\begin{equation}
 AB=BA=0,
 \qquad
 \rank A^2=\rank A,
\end{equation}
则
\begin{equation}
 \rank(A+B)=\rank A+\rank B.
\end{equation}
\end{theorem}

\begin{proof}
条件 $\rank A^2=\rank A$ 等价于
\begin{equation}
 \Image A\cap\Ker A=\{0\}.
\end{equation}
由 $AB=0$ 得 $\Image B\subseteq\Ker A$，所以
\begin{equation}
 \Image A\cap\Image B=\{0\}.
\end{equation}
对转置矩阵使用 $BA=0$ 与 $\rank(A^T)^2=\rank A^T$，同样得到行空间之交为零。由矩阵和秩等号的判别即得结论。
\end{proof}

\subsection{LU 分解}

\begin{theorem}
设 $A\in\F^{n\times n}$ 的各阶顺序主子式均非零：
\begin{equation}
 \det A_{1:k,1:k}\neq0,
 \qquad k=1,\ldots,n.
\end{equation}
则存在唯一分解
\begin{equation}
 A=LU,
\end{equation}
其中 $L$ 是对角元全为 $1$ 的下三角矩阵，$U$ 是可逆上三角矩阵。
\end{theorem}

\begin{proof}
顺序主子式非零保证 Gaussian 消元无需换行，所有消元乘子组成 $L$，消元后的阶梯矩阵为 $U$。若另有 $A=\widetilde L\widetilde U$，则
\begin{equation}
 \widetilde L^{-1}L=\widetilde U U^{-1}.
\end{equation}
左端是单位下三角，右端是上三角，只能都等于 $I$。
\end{proof}

\begin{remark}
若允许置换，一般可写成
\begin{equation}
 PA=LU.
\end{equation}
数值计算中选主元的目的不仅是避免零主元，也用于控制舍入误差。
\end{remark}

\subsection{Householder 变换与 QR 分解}

\begin{definition}
对单位向量 $w\in\R^n$，矩阵
\begin{equation}
 H=I-2ww^T
\end{equation}
称为 Householder 反射。
\end{definition}

\begin{proposition}
Householder 矩阵满足
\begin{equation}
 H^T=H,
 \qquad
 H^2=I,
 \qquad
 H^TH=I.
\end{equation}
它把 $w$ 方向取反，并保持 $w^\perp$ 中的向量不变。
\end{proposition}

\begin{theorem}
给定非零向量 $x\in\R^n$，可选单位向量 $w$ 使相应 Householder 反射满足
\begin{equation}
 Hx=\pm\norm{x}_2e_1.
\end{equation}
例如取
\begin{equation}
 v=x+\operatorname{sgn}(x_1)\norm{x}_2e_1,
 \qquad
 w=\frac{v}{\norm v_2},
\end{equation}
即可避免严重的相消误差。
\end{theorem}

\begin{remark}
依次对矩阵第一列、第二列以下部分等应用 Householder 反射，可将 $A$ 化为上三角或上梯形矩阵：
\begin{equation}
 H_k\cdots H_1A=R.
\end{equation}
于是
\begin{equation}
 A=QR,
 \qquad
 Q=H_1\cdots H_k
\end{equation}
是 QR 分解。该方法比经典 Gram--Schmidt 通常更稳定。
\end{remark}
